% Generated draft: requires mathematical and source review.
\documentclass[11pt]{article}
\usepackage[margin=25mm]{geometry}
\usepackage{fontspec}
\setmainfont{Latin Modern Roman}
\usepackage{amsmath,amssymb,mathtools,unicode-math}
\setmathfont{Latin Modern Math}
\usepackage{enumitem,needspace,xurl,hyperref}
\setlength{\parindent}{0pt}
\setlength{\parskip}{5pt}
\setlength{\emergencystretch}{3em}
\newcommand{\N}{\mathbb{N}}
\newcommand{\Z}{\mathbb{Z}}
\newcommand{\Q}{\mathbb{Q}}
\newcommand{\R}{\mathbb{R}}
\newcommand{\C}{\mathbb{C}}
\newcommand{\F}{\mathbb{F}}
\newcommand{\A}{\mathbb{A}}
\newcommand{\PP}{\mathbb{P}}
\newcommand{\E}{\mathbb{E}}
\newcommand{\eps}{\varepsilon}
\newcommand{\dd}{\,\mathrm{d}}
\begin{document}
\section*{UnsolvedMath ID: 30005201}
\textbf{Source problem code:} catalogue entry
\subsection{Definitions and mathematical statement}
The supplied record does not identify a recoverable open problem. It contains a truncated conjecture citation followed by statements labelled Theorems 4 and 5. The following records their mathematical context and a qualified reconstruction; it does not assert that either theorem is an open conjecture.

An \(M\)-\(M\) bimodule is a Hilbert space with commuting normal left and right actions of a von Neumann algebra \(M\); the right action is equivalently an action of \(M^{\mathrm{op}}\). For a closed Hilbert subspace \(H_1\subseteq H_2\), write \(H_2\ominus H_1=H_2\cap H_1^\perp\). A von Neumann subalgebra is maximal amenable if it is amenable and is contained in no strictly larger amenable von Neumann subalgebra of the ambient algebra.

The source uses \(L(\mathbb F_t)\) for \(t>1\), apparently referring to interpolated free group factors, but does not define this notation or explicitly specify the parameter domain beyond \(t>1\). Retaining its notation, put
\[
 M=L(\mathbb F_t),\qquad
 \mathcal H_P=L^2(M)\ominus L^2(P),\qquad
 \mathcal C_P=L^2(P)\otimes L^2(P).
\]
Here the \(L^2\)-spaces are those associated with the trace understood in the source. The coarse \(P\)-\(P\) bimodule structure on \(\mathcal C_P\) is
\[
 a\cdot(\xi\otimes\eta)\cdot b=(a\xi)\otimes(\eta b),
 \qquad a,b\in P.
\]
A plausible reconstruction of the corrupted display labelled Theorem 4 is: for every \(t>1\) and every maximal amenable subalgebra \(P\subseteq L(\mathbb F_t)\), there exists a \(P\)-bimodular isometric linear embedding
\[
 U:\mathcal H_P\longrightarrow\bigoplus_{n=1}^{\infty}\mathcal C_P,
 \qquad U(a\xi b)=aU(\xi)b.
\]
This is only a conditional reconstruction. The supplied prose defines its comparison symbol using a bimodular \emph{unitary} map \(H\to K\). With the usual surjectivity convention for a unitary, that would instead require a bimodule isomorphism onto the entire target. The damaged display and this terminology do not determine which assertion was intended. The source calls the property coarseness and attributes this terminology to its unidentified reference [20], with an informal interpretation as a maximally random position of the subalgebra.

For the abelian formulation, let \((M,\tau)\) be a tracial von Neumann algebra and let \(A\subseteq M\) be a maximal abelian \(*\)-subalgebra. As in the source, choose a presentation \(A=L^\infty(X,\mu)\), with \(X\) compact Hausdorff and \(\mu\) a Borel probability measure. On \(\mathcal H_A=L^2(M,\tau)\ominus L^2(A,\tau)\), consider
\[
 \pi(f\otimes g)\xi=f\xi g,\qquad f,g\in C(X).
\]
Write \(E\) for the associated projection-valued spectral measure on \(X\times X\). A left/right measure is a Borel probability measure \(\nu\) on \(X\times X\) with the same null sets as \(E\): for every Borel set \(B\),
\[
 \nu(B)=0\quad\Longleftrightarrow\quad E(B)=0.
\]
Thus the term denotes a measure class, and any representative may be used. The background also asserts that the marginals of \(E\) have the same null sets as \(\mu\); the hypotheses needed for that general assertion are not established by the supplied excerpt.

The statement labelled Theorem 5 is: for every \(t>1\), if \(M=L(\mathbb F_t)\) and \(A\subseteq M\) is abelian and maximal amenable, then, for any presentation \(A=L^\infty(X,\mu)\) as above with \(X\) additionally compact metrizable, every left/right measure \(\nu\) satisfies
\[
 \nu\ll\mu\otimes\mu,
\]
meaning that \((\mu\otimes\mu)(B)=0\) implies \(\nu(B)=0\) for every Borel \(B\subseteq X\times X\). The source introduces left/right measures for maximal abelian subalgebras, whereas this theorem explicitly assumes abelianness and maximal amenability; their connection is left implicit in the excerpt.

No separate unresolved conjecture can be specified from the leading fragment ``Conjecture 1.12] and [20, Conjecture 5.2]''. The supplied open label conflicts with supplied verification reporting a positive resolution of the coarseness conjecture.
\subsection{Short English statement}
The excerpt reports coarseness and absolute continuity of left/right measures for maximal amenable subalgebras of \(L(\mathbb F_t)\), with \(t>1\). Its coarseness formula is damaged, and its conjecture references are incomplete. The supplied record labels the target open but also reports that the conjecture was resolved. A unique open problem cannot be recovered from this evidence.
\subsection{Sources}
\begin{itemize}
\item Local UnsolvedMath registry, supplied complete record: permanent ID \(30005201\), exact source problem code \texttt{catalogue entry}, title ``Resolved Coarseness Conjecture for Free Group Factors''. No proposer is supplied; the proposed year is recorded as 2022.
\item Supplied report citation: \(C^*\)-Algebras (2023), \emph{Oberwolfach Reports}, pp.\ 2059--2127, DOI \url{https://doi.org/10.4171/owr/2022/36}. This URL is also the supplied literature source labelled ``original Oberwolfach report''. The supplied report title/year is retained; the DOI contains 2022.
\item Supplied literature assessment, checked exactly \texttt{2026-08-22}: the extraction conflates conjecture references with theorem statements and does not isolate a verifiable mathematical claim; its triage therefore labels the unresolved target open. This is supplied dated curation, not an independently verified assessment. The background explicitly requests independent verification of time-sensitive status and cited claims.
\item Supplied statement-verification note: the official report is said to state that the coarseness conjecture was resolved positively and that the extracted passage is not an open problem. This report about the source has not been independently verified here. The record marks its statement status as \texttt{unrecoverable}, while retaining status \texttt{open}.
\item The excerpt refers to ``Conjecture 1.12'', ``[20, Conjecture 5.2]'', Theorem 4, and Theorem 5. It supplies no bibliography identifying [20], no complete citation for Conjecture 1.12, and no individual attribution for these statements.
\end{itemize}
\end{document}
